\section{Methodology}

Our experiments proceed in four stages. We first reconstruct the value axis following the method of \citet{jiang2026value}, freezing the axis once it has been verified. We then test whether the frozen axis transfers to long-horizon agentic trajectories, first at the final step of each trajectory and then across all trajectory positions. Next, we compare the resulting internal projection with the agent's own verbalized probability of success, testing whether the verbalized estimate aligns with or diverges from the internal signal. Finally, we fit linear probes directly on the same activations for comparison.


\paragraph{Reconstructing the value axis.} We reconstruct the value axis on Qwen-3-8B following the difference-in-means construction of \citet{jiang2026value}. We create our own in-context RL dataset using Claude Opus 4.6 and verify the reconstruction on a held-out test set. We require the highest AUROC to be at least 0.90, with peak performance occurring in the middle to late layers, consistent with the findings of \citet{jiang2026value}. The layer with the highest AUROC in our reconstruction is designated as the primary layer for all subsequent analyses. Results at other layers are reported to study how the signal performs throughout the model. Once the axis is verified, we do not refit it at any later stage.

\paragraph{Transfer to agentic trajectories.} Following the evaluation procedure of \citet{jiang2026value}, we define the alignment score of trajectory step $t$ as the mean cosine similarity between the frozen axis $v^{(\ell)}$ and the layer-$\ell$ activations of the tokens generated by the agent at that step:

\begin{equation}
A(t) = \frac{1}{|G_t|} \sum_{i \in G_t} \cos\!\left(h^{(\ell)}_i,\; v^{(\ell)}\right),
\end{equation}

where $G_t$ is the set of agent-generated token positions in step $t$. Higher values of $A(t)$ indicate stronger alignment between the agent's internal activations and the value axis. We restrict this mean to agent-generated tokens, excluding tool output and scaffold text. For comparison, we also report the alignment score at the final generated token alone. We first evaluate whether $A(T)$ at the final step T separates resolved from unresolved trajectories, as this is the most favorable case for detecting the signal. We then compute the same metric at every relative position, binned into five equally sized position bins. A failure to separate outcomes at the final step is not immediately interpreted as evidence of failed transfer unless separation is also absent across the rest of the trajectory.

\paragraph{Comparison against verbalized confidence.} 


After the generation and completion of all trajectories, we go back and elicit a post hoc verbalized estimate of how likely the model was to solve the task. 
For each trajectory prefix, ending at step t, we prompt the model in a separate forward pass
to state a single number in $[0, 1]$ corresponding to its probability of eventually resolving the task (prompt in Appendix). By eliciting the model's verbalized success estimate post hoc, we leave the original trajectory clean and eliminate the risk of producing and then studying a contaminated run. But by doing so, we introduce an asymmetry worth stating plainly: the internal projection is read from the rollout pass itself, whereas the verbalized estimate comes from a separate pass over the same prefix with an elicitation prompt appended. Both are nonetheless scored against the same ground-truth outcomes at each position, testing whether the internal state carries outcome information beyond what the agent reports. 



\paragraph{Fitted-probe reference.} 

To test in-domain linear decodability, we fit a supervised linear probe directly on activations collected from the SWE-bench trajectories. This asks whether the model's hidden states contain linearly accessible information about eventual task resolution, independently of whether that information is captured by the transferred value axis. For trajectory $j$ at step $t$, we represent the agent's internal state using the mean layer-$\ell$ activation over agent-generated tokens, $\bar{h}_{j,t}^{(\ell)}$. We fit a logistic-regression probe of the form

\begin{equation}
\hat{p}_{j,t}^{(\ell)} = \sigma\!\left(w^{(\ell)\top}\bar{h}_{j,t}^{(\ell)} + b^{(\ell)}\right),
\end{equation}

where $w^{(\ell)}$ and $b^{(\ell)}$ are learned probe parameters, $\sigma$ is the logistic sigmoid function, and $\hat{p}_{j,t}^{(\ell)}$ is the predicted probability that trajectory $j$ will eventually resolve the task. Each activation is assigned the final trajectory label $y_j \in \{0,1\}$, where $y_j=1$ denotes a resolved trajectory and $y_j=0$ denotes an unresolved trajectory.

We divide the 60 SWE-bench tasks equally between training and test partitions. The split is performed at the task level, where all steps and all five random-seed trajectories associated with a task are assigned to the same partition. The probe is fitted exclusively on the 30 training tasks and evaluated on the 30 held-out tasks. We report held-out AUROC against the final SWE-bench outcomes using the same trajectory positions and layers used to evaluate the frozen value-axis projection.

% TODO if applicable: Insert details about standardization, regularization, and layer selection.

The fitted probe and frozen value axis answer complementary questions. If both separate signals resolved from unresolved trajectories, outcome information is linearly decodable and the previously constructed value axis transfers to the agentic setting. If only the fitted probe succeeds, outcome information remains linearly decodable but is not captured by the transferred value direction. If only the frozen value axis succeeds, this supports representational transfer but indicates that the supervised probe did not reliably recover the signal from the available in-domain training data. If neither succeeds, we find no reliable linearly decodable outcome signal under the evaluated activations and experimental setup.

\section{Experimental Setup}

\paragraph{Model and scaffold.} Our experiments use Qwen-3-8B, the same model on which the value axis was originally constructed. This removes differing model architecture as a confounding factor. Trajectories are generated by Qwen-3-8B running in the mini-SWE-agent harness, version 2.4.5. We set a step limit of 60, consistent with evidence that the agent only continues unsuccessfully beyond this point \citep{agentstop2026}.We require that all trajectories are generated live by Qwen-3-8B itself.


\paragraph{Tasks and trajectories.} We select 60 tasks from SWE-bench Verified \citep{jimenez2024swebench}, distributed based on human-annotated time-to-solve estimates. We select 25 easy (<15 minutes), and 35 medium (15 minutes--1 hour), tasks,  excluding the tasks annotated at 1--4 hours and >4 hours. Qwen-3-8B resolves only a minority of SWE-bench Verified instances, and in the harder tiers, its expected resolution rate approaches zero. Because of this, a difficulty group where there is no success class contributes nothing to outcome separation and nothing to within-task findings. Also, if difficulty correlates with outcome, a probe can separate outcomes by detecting difficulty from the prompt alone, without tracking the shape of the trajectory at all. Thus, since tasks in the hardest difficulty tiers are substantially less likely to produce both outcomes, they are less useful for the analyses in this work, and it is for this reason that they are excluded. Within each difficulty group, we sample across repositories with a per-repository cap of 8 tasks. Selection is based on the human-annotated difficulties and no model outputs, activations, or trajectories were examined to select the tasks. We generate 5 independent trajectories per task for 300 generated trajectories; the final analyzed trajectory count after exclusions is reported in Section [REF]. The five trajectories per task hold everything constant except the random seed. Sampling uses temperature 0.6, top-$p$ 0.95, and top-$k$ 20, with the task, repository, and scaffold held fixed. Outcomes are binary, labeled as resolved or unresolved, determined by SWE-bench's FAIL\_TO\_PASS and PASS\_TO\_PASS test criteria. Repeated trajectories per task allow for within-task contrasts, where the same task is resolved on one seed and not another. Resolved and unresolved trajectories are compared while holding task difficulty fixed to test whether the signal tracks trajectory shape rather than static task difficulty.


\paragraph{Statistical reporting.} Because resolved trajectories are the minority class, all separation is reported as AUROC against a baseline that always predicts unresolved. To estimate the uncertainty of our results, we use bootstrapping. Because within-task trajectories are similar, we resample the tasks and recompute AUROC repeatedly. We then report 95\% bias-corrected and accelerated (BCa) confidence intervals on all reported AUROCs. Additionally, we perform a permutation test, randomly reassigning the resolved and unresolved labels while keeping the trajectories together. This helps ensure that our results capture the desired signal and not random chance.